\section{Lethai−ar — ``The Crimson Valley''}
łabel{chap:lethai−ar}

\textbf{Key Demographic: } Lethai−ar
\textbf{Capital City: } Ananse

\subsection*{Quick Reference}
\textbf{Tagline:} \textit{The bridge remembers every step. The silk knows every promise.}

\begin{quote}
  \textbf{Silk−Warden (Elite):}  
  ``The web−law is not written on parchment. It is woven into every bridge, every strand, every knot. When a promise is made, the valley remembers. When a promise is broken, the silk tightens. We are not judges. We are the loom. The spider spins; we only follow the thread.''
\end{quote}

\begin{quote}
  \textbf{Mark−Cutter (Commoner):}  
  ``Each line on my skin is a promise I witnessed. Each line is a story I carry. When I die, they will cut these marks from my body and bury them at the Ukwe's roots. Then the tree will remember what I witnessed. I hope it is enough. The thread does not break; it only passes to the next hand.''
\end{quote}

\textbf{Genre / Mood:} High−fantasy legal drama, spidersilk noir, ancestral witness.

\textbf{Starting Location:} The central silk bridge, where the gorge is deepest and the wind carries the smell of rust. A Silk−Warden blocks your path, her mask painted with roots: ``The web−law demands a toll. Not in coin – in a promise. Swear that you will not use this valley to carry steel to Ashaan, or turn back now. The thread is watching. So is the spider.''

\vspace{2mm}
\begin{tcolorbox}[colback=black!3,colframe=blue!60!black,title={Theme & Atmosphere}]
The Crimson Valley is a deep jungle rift between Sekogo and Oshiira – a wound of red stone, carved by a river that runs rust‑coloured after every rain. Here, the Lethai−ar guard their borders with fierce neutrality, their chestnut‑brown skin blending with the canyon walls. Silk bridges span the gorges, woven from spider‑sap and oath‑hair, strong enough to hold a caravan and light enough to sway in the monsoon wind. The web−law is unwritten, carried in the memory of each warden; it keeps both empires from using the valley as a shortcut to war.

Unlike their northern kin who fell from grace, the Lethai−ar are adopted children of the old powers – not fallen, but chosen to carry memory. Each can name the specific mountain, river, or grove where their ancestor left the spirit world to become mortal. Their skin‑tattoos (Marks) are woven with red ochre and jungle sap, each line a promise witnessed, each curve a story recorded. When a Lethai−ar dies, their Marks are cut from the skin and buried at the Ukwe's roots, so the tree may remember what they witnessed.

The valley itself is alive. The silk remembers every footfall; the river whispers the names of those who have crossed; and in the oldest bridges, the spider – Inaea, the weaver of bonds – still spins. The Lethai−ar do not worship her, but they honour the thread. They say that every promise is a strand in her web, and that a broken oath sends a tremor through the whole valley.

\vspace{2mm}
\textbf{What you notice first:} The whisper of silk ropes in the wind. The smell of red ochre and damp stone. The way every Lethai−ar touches their left wrist when they speak a name – a gesture of witness.

\textbf{The one rule that kills newcomers:} Never step onto a silk bridge without first naming a witness. The bridge will remember your footsteps. If you fall, it will ask why no one was watching.
\end{tcolorbox}

\subsection*{Regional Mood: The Weight of Witness}

Power here flows from witnessed promises, silk strands, and the web−law's quiet authority. The valley is neutral ground – Oshiiran grain barges do not float here, and Ashaani armies dare not march. The Lethai−ar trade silk for honey with Sekogo druids and stonesong for poison with Ngomebe masons. Their Marks are not decoration; they are a public memory, written on skin, read by anyone who knows the phase of the crimson moon. A Lethai−ar who bears a Mark they cannot fulfil loses standing; a warden who cuts a bridge without cause is exiled. But the true authority is the web itself – the invisible threads that connect every oath ever sworn in the valley. The Lethai−ar do not create the web. They tend it.

\subsection*{Prominent NPCs of the Crimson Valley}
łabel{sec:lethaiar−npcs}

\begin{tabular}{|p{0.2\textwidth}|p{0.25\textwidth}|p{0.3\textwidth}|p{0.15\textwidth}|}
ℎline
\textbf{Name} & \textbf{Role/Title} & \textbf{Motivation} & \textbf{Rumoured Quirk} \\
ℎline
Warden Anara & Silk‑warden of the central span & Keep the Oshiiran grain barges from crossing & Her left arm is covered in Marks from wrist to shoulder. She will not say what they mean. \\
ℎline
Mark‑Cutter Tebe & Tattooist who records promises & Witness a treaty between Sekogo and Ngomebe & She uses a needle of silk‑fiber and ink made from crushed crimson stone. The ink glows under the full moon. \\
ℎline
Bridge‑Keeper Solan & Guardian of the eastern crossing & Prevent Ashaani scouts from using the valley & He has not slept in nine days. He says the bridge whispers when someone lies. \\
ℎline
Elder Mira & Memory‑keeper of the valley's spirits & Record the name of every ancestor who walked the spirit path & She carries a bundle of silk cords, each tied to a different knot. The knots are the names. \\
ℎline
The Spider's Child & Tulkani trader of oath‑silk (whisperer of Inaea) & Buy a bolt of crimson silk for a story – and to listen & She is not Lethai−ar but has lived in the valley for thirty years. The wardens tolerate her. She never closes her eyes. \\
ℎline
The Unmarked & A Lethai−ar who refused any Marks & An outcast who lives in the gorge below the bridges & Her skin is blank. She is invisible to the web‑law – and to the valley's spirits. She carries a single knotted thread. \\
ℎline
\end{tabular}

\subsection*{Names of the Crimson Valley}
łabel{sec:lethaiar−names}

Lethai−ar names are soft, with rolling consonants and vowel endings. Surnames often denote a bridge, a gorge, or a spirit‑ancestor. Epithets refer to Marks or witnessed deeds.

\begin{tabular}{|p{0.25\textwidth}|p{0.25\textwidth}|p{0.2\textwidth}|p{0.2\textwidth}|}
ℎline
\textbf{First Names (Male)} & \textbf{First Names (Female)} & \textbf{Surnames} & \textbf{Epithets} \\
ℎline
Solan & Anara & Silkweave & \textit{the Unbroken} \\
ℎline
Tebe & Mira & Gorge‑born & \textit{the Ninth Witness} \\
ℎline
Liam & Nira & Root‑memory & \textit{the Marked} \\
ℎline
Veth & Liana & Spider‑kin & \textit{the Bridge‑Heart} \\
ℎline
Renn & Sora & Crimson‑touch & \textit{the Unweaver} \\
ℎline
Thorn & Aria & Oath‑strand & \textit{the Still Watcher} \\
ℎline
\end{tabular}

\paragraph*{(Bridge/Gorge/Roots)} Silk bridge spanning a crimson gorge – ropes sing in the wind; a Lethai−ar village built into the canyon wall; the Ukwe's roots clutching the water below the central span.

\section*{Spades — Places}
łabel{sec:lethaiar−places}
\begin{enumerate}
\setcounter{enumi}{1}
ℹtem \textbf{Central Silk Bridge} – A wide span woven from crimson and grey strands. The bridge hums when a promise is made on it.
  ℎfill \textit{The pitch tells the warden if the promise is true. If the thread vibrates at the ninth frequency, the spider is listening.} \texttt{[WITNESS]}
ℹtem \textbf{Mark‑Cutter's Studio} – A cave with a stone table and pots of ochre.
  ℎfill \textit{The walls are carved with Marks that were never cut into skin. They are stories the cutter refused to witness. The stone is warm to the touch.} \texttt{[CRAFT]}
ℹtem \textbf{Ukwe's Riverside Root} – A massive root that dips into the crimson river.
  ℎfill \textit{The root is warm. Touch it, and the tree will show you one forgotten promise from a dead Lethai−ar. The water ripples in response.} \texttt{[MEMORY]}
ℹtem \textbf{Abandoned Silk Span} – A broken bridge hanging over the gorge. The cut ends are frayed.
  ℎfill \textit{The cut was not clean – someone broke an oath mid‑crossing. The strands still twitch.} \texttt{[DEATH]}
ℹtem \textbf{Spirit‑Anchor Stone} – A boulder where the first ancestor left the spirit world.
  ℎfill \textit{The stone is smooth. It is warm even in winter. Lethai−ar leave offerings of red ochre here – and sometimes a single knotted thread.} \texttt{[RITUAL]}
ℹtem \textbf{Warden's Watchtower} – A stone needle built into the canyon wall, with a signal mirror.
  ℎfill \textit{The mirror flashes in code. The code is the web‑law. No outsider has broken it. The code changes when the spider dreams.} \texttt{[AUTHORITY]}
ℹtem \textbf{Hollow of the Unmarked} – A cave below the eastern bridge, hidden by a waterfall.
  ℎfill \textit{The cave has no silk. The Unmarked lives here, weaving nets from her own hair. Her hair is knotted.} \texttt{[MYSTERY]}
ℹtem \textbf{Crimson River Ford} – A shallow crossing where the water is clear, not red.
  ℎfill \textit{The ford is sacred. No oath sworn here can be broken. The water would turn black – and the spider would feel it.} \texttt{[OATH]}
ℹtem \textbf{Elder's Cord‑House} – A longhouse where the memory‑knots are stored.
  ℎfill \textit{Each cord is a name. The cords are hung from the ceiling. They move in no wind. When a name is forgotten, a cord falls.} \texttt{[MEMORY]}
ℹtem[\card{♠}{J}] \textbf{The Spider's Stall} – A Tulkani merchant's tent at the valley's southern exit. She sells silk to anyone.
  ℎfill \textit{The price is a story. The story must be about a thread that held – or a thread that broke. She never blinks.} \texttt{[TRADE]}
ℹtem[\card{♠}{Q}] \textbf{Bridge of the Ninth Strand} – A legend: a bridge woven from a single hair.
  ℎfill \textit{It appears only at dawn on the first day of spring. Those who cross it gain a new name – and a new thread to their web.} \texttt{[LEGEND]}
ℹtem[\card{♠}{K}] \textbf{The Cut Stone} – A boulder split by an oath‑breaker's fall.
  ℎfill \textit{The split is shaped like a human body. The stone weeps water when the oath‑breaker's name is spoken. The water tastes of rust.} \texttt{[DEATH]}
ℹtem[\card{♠}{A}] \textbf{The First Bridge (Myth)} – Where the first Lethai−ar wove the first strand.
  ℎfill \textit{It is invisible. Only those who have witnessed nine true promises can see it. The spider sits at the far end, waiting.} \texttt{[MYTH]}
\end{enumerate}

\paragraph*{(Warden/Cutter/Trader)} Silk‑warden with a mask painted with the Ukwe's roots; Mark‑cutter with a needle of silk‑fiber; Tulkani spider‑trader with a bundle of oath‑cords.

\section*{Hearts — People & Factions}
łabel{sec:lethaiar−people}
\begin{enumerate}
\setcounter{enumi}{1}
ℹtem \textbf{Silk−Warden} – Mask carved from crimson wood, painted with oath‑names.
  ℎfill \textit{She never blinks when standing on the bridge. The bridge does not forgive distraction. She counts her breaths in eights.} \texttt{[AUTHORITY]}
ℹtem \textbf{Mark−Cutter} – Hands stained red with ochre, eyes that measure skin.
  ℎfill \textit{She asks each client: ``What promise do you want the world to remember?'' The answer becomes a Mark – and a thread in the web.} \texttt{[CRAFT]}
ℹtem \textbf{Bridge−Keeper} – Old, grey‑haired, missing three fingers from a snapped strand.
  ℎfill \textit{He carries a whistle made from a spider's fang. The whistle calls the bridge to tighten – or to loosen. He never uses it.} \texttt{[WITNESS]}
ℹtem \textbf{Elder (Memory‑Keeper)} – Bundles of silk cords hang from her belt.
  ℎfill \textit{She knows every name carved into every root of the Ukwe. She has forgotten her own. The cords remember it for her.} \texttt{[MEMORY]}
ℹtem \textbf{Tulkani Spider (Trader)} – Wears a shawl of crimson silk, speaks in riddles.
  ℎfill \textit{She pays for silk with stories. The stories become new strands. The bridge grows. She never closes her eyes – she says the spider does not sleep.} \texttt{[TRADE]}
ℹtem \textbf{The Unmarked} – A Lethai−ar without Marks, living in the hollow.
  ℎfill \textit{She whispers to the waterfall. The waterfall whispers back. No one knows what they say. Her thread is her own.} \texttt{[MYSTERY]}
ℹtem \textbf{Youth Seeking a Mark} – A young Lethai−ar on their first witness journey.
  ℎfill \textit{They carry a single strand of undyed silk. When they witness a promise, they will dye it crimson. The dye is their own blood.} \texttt{[COMMUNITY]}
ℹtem \textbf{Oath‑Breaker's Ghost} – A translucent figure pacing the broken bridge.
  ℎfill \textit{It repeats the words of the promise it broke. The words are getting quieter. When they stop, the thread will snap.} \texttt{[DEATH]}
ℹtem \textbf{Sekogo Honey Trader} – A druid with a jar of dark honey, trading for silk poison.
  ℎfill \textit{He offers the honey as a gift. The wardens test it by feeding a drop to the bridge. The bridge sways if it is poisoned – or if the gift is false.} \texttt{[TRADE]}
ℹtem[\card{♥}{J}] \textbf{Ngomebe Stone‑Singer} – A mason with a tuning fork, here to buy silk for sledge‑wax seals.
  ℎfill \textit{He taps the bridge with his fork. The note tells him how many promises it has witnessed. The note is always an octave below the spider's hum.} \texttt{[CRAFT]}
ℹtem[\card{♥}{Q}] \textbf{The First Ancestor's Echo} – A voice heard at the Spirit‑Anchor Stone at midnight.
  ℎfill \textit{It asks: ``Why did you leave the spirit world?'' The answer must be a promise. The spider listens.} \texttt{[RITUAL]}
ℹtem[\card{♥}{K}] \textbf{The Spider's Daughter} – A half‑Lethai−ar child raised by the Tkanki trader. She has no Marks.
  ℎfill \textit{Her skin is blank. She is learning to read knots – and to spin her own thread. She never asks twice.} \texttt{[COMMUNITY]}
ℹtem[\card{♥}{A}] \textbf{The Ninth Warden (Legend)} – A warden who cut her own bridge to stop an invasion.
  ℎfill \textit{She is said to walk the valley at dusk, carrying a coil of silk. She offers one free crossing – but the crossing leads to the spider's web.} \texttt{[LEGEND]}
\end{enumerate}

\paragraph*{(Strand/Challenge/Fall)} A silk strand snaps – lose 1 Supply to repair; a Lethai−ar challenges you to recite the web‑law; a traveler falls from the bridge – the party must witness their last words.

\section*{Clubs — Complications/Threats}
łabel{sec:lethaiar−complications}
\begin{enumerate}
\setcounter{enumi}{1}
ℹtem \textbf{Snapped Strand} – A rope on a bridge breaks. Lose 1 Supply to repair.
  ℎfill \textit{The break is clean. Someone cut it with a blade. The blade was from Ashaan – or from a broken promise.} \texttt{[DISASTER]}
ℹtem \textbf{Web−Law Challenge} – A warden blocks the bridge: ``Recite the ninth clause.''
  ℎfill \textit{The ninth clause is unwritten. Each warden carries a different version. You must guess. The spider knows the true one.} \texttt{[SOCIAL]} \texttt{[RIDDLE]}
ℹtem \textbf{Fall from the Bridge} – A traveler slips. The bridge asks: ``Who witnessed this?''
  ℎfill \textit{If no one answers, the traveler's name is erased from the Ukwe's roots – and the spider weaves a new strand.} \texttt{[DEATH]} \texttt{[WITNESS]}
ℹtem \textbf{Oshiiran Surveyors} – Grain auditors arrive with a map, claiming the valley's water rights.
  ℎfill \textit{They carry a treaty stamped with the Widow's seal. The seal is a forgery. The bridge hums in warning.} \texttt{[AUTHORITY]} \texttt{[POLITICS]}
ℹtem \textbf{Ashaani Scouts (Disguised)} – Veiled Sister agents posing as spice merchants.
  ℎfill \textit{They do not touch their left wrist when speaking a name. The wardens notice. The thread tightens.} \texttt{[STEALTH]} \texttt{[POLITICS]}
ℹtem \textbf{Mark−Cutter's Witness Called} – A Lethai−ar's Marks begin to glow red. A promise is overdue.
  ℎfill \textit{The witness is not present. The party is the closest witness. The duty transfers – as does the thread.} \texttt{[WITNESS]} \texttt{[COMMUNITY]}
ℹtem \textbf{Bridge Poisoning} – Someone smears a solvent on the silk strands. The bridge begins to fray.
  ℎfill \textit{The solvent is made from Sekogo honey and Ashaani dust. The Spider knows who sold it. The threads are weeping.} \texttt{[MAGIC]} \texttt{[DISASTER]}
ℹtem \textbf{Elder's Memory Fails} – The cord‑keeper forgets a name. The knot unravels.
  ℎfill \textit{The forgotten name is an old promise. The Ukwe's root trembles. Something is coming – or someone.} \texttt{[MEMORY]} \texttt{[DEATH]}
ℹtem \textbf{The Unmarked's Bargain} – The outcast offers to guide the party through the gorge. Her price is a promise.
  ℎfill \textit{The promise: ``Do not speak my name after I die.'' Refuse, and she leaves you to the spiders. Accept, and she ties a knot.} \texttt{[SOCIAL]} \texttt{[WITNESS]}
ℹtem[\card{♣}{J}] \textbf{Crimson River Floods} – Rain turns the ford to a torrent. The bridge is the only crossing.
  ℎfill \textit{The water is red. Faces appear in the foam. They are oath‑breakers, still trying to cross. The spider watches.} \texttt{[WEATHER]} \texttt{[DEATH]}
ℹtem[\card{♣}{Q}] \textbf{Warden Schism} – Two wardens disagree on the web‑law. They cut each other's bridges.
  ℎfill \textit{The valley is now divided. The party must choose which side to cross – or stay stuck forever. The spider waits.} \texttt{[AUTHORITY]} \texttt{[POLITICS]}
ℹtem[\card{♣}{K}] \textbf{The First Bridge Appears} – The mythical bridge manifests at dawn. Crossing it gives a new name.
  ℎfill \textit{The new name is a promise. The promise is to the First Ancestor. The ancestor will collect – and the spider will weave.} \texttt{[MYTH]} \texttt{[MAGIC]}
ℹtem[\card{♣}{A}] \textbf{The Ninth Warden Returns} – The legend appears, offering a free crossing.
  ℎfill \textit{The crossing leads to the spider's web. The Hollow asks: ``What promise did you break?'' Answer poorly, and you become a new strand.} \texttt{[DEATH]} \texttt{[LEGEND]}
\end{enumerate}

\paragraph*{(Cord/Mask/Shard)} Silk truth‑cord (enforce any oath); warden's mask (safe passage on any bridge for one crossing); Ukwe root shard (one question about a forgotten promise).

\section*{Diamonds — Rewards/Leverage}
łabel{sec:lethaiar−rewards}
\begin{enumerate}
\setcounter{enumi}{1}
ℹtem \textbf{Silk Truth‑Cord} – A braided strand dyed crimson. Use it to enforce any oath.
  ℎfill \textit{The cord tightens when the oath is broken. It will not stop until the promise is kept – or the spider releases it.} \texttt{[OATH]}
ℹtem \textbf{Warden's Mask} – A wooden mask painted with roots. Safe passage on any bridge for one crossing.
  ℎfill \textit{The mask has no eye holes. You must trust the bridge to see for you. The bridge sees through the spider's eyes.} \texttt{[TRAVEL]}
ℹtem \textbf{Ukwe Root Shard} – A piece of the tree's root that grows near the river.
  ℎfill \textit{Touch it, and the Ukwe will answer one question about a forgotten promise. The answer is always true – and always a thread.} \texttt{[MEMORY]}
ℹtem \textbf{Mark‑Cutter's Needle} – A silk‑fiber needle. Use it to record a witnessed promise as a temporary Mark.
  ℎfill \textit{The Mark fades after one season. The story remains – and the spider remembers.} \texttt{[CRAFT]}
ℹtem \textbf{Bridge‑Keeper's Whistle} – A fang whistle. Blow it to tighten any silk bridge for one scene.
  ℎfill \textit{The whistle works once. The bridge will remember the favour – and the spider will know.} \texttt{[AUTHORITY]}
ℹtem \textbf{Elder's Memory Cord} – A knotted silk cord that contains a forgotten name.
  ℎfill \textit{Untie the knot to learn the name. The name is a promise. Fulfil it, or the cord will retie itself – and the spider will watch.} \texttt{[MEMORY]}
ℹtem \textbf{Spider's Story‑Strand} – A bolt of crimson silk bought with a story.
  ℎfill \textit{The silk is soft. It glows under moonlight. Trade it for safe passage through Tkanki camps – or give it to the spider.} \texttt{[TRADE]}
ℹtem \textbf{The Unmarked's Net} – A net woven from human hair.
  ℎfill \textit{Throw it to catch a spirit. The net can hold a bound spirit for one hour. The spirit will bargain for freedom – with a promise.} \texttt{[MAGIC]}
ℹtem \textbf{Crimson Ochre Pot} – A clay pot of red pigment. Use it to Mark a promise on stone or skin.
  ℎfill \textit{The ochre never dries. The promise remains witnessable forever – until the spider cuts the thread.} \texttt{[CRAFT]}
ℹtem[\card{♦}{J}] \textbf{Ninth Strand Fragment} – A piece of the legendary bridge. Once per campaign, cross any gap as if on solid ground.
  ℎfill \textit{The fragment crumbles after use. The first bridge will remember you – and the spider will ask for a toll.} \texttt{[LEGEND]}
ℹtem[\card{♦}{Q}] \textbf{The Cut Stone's Water} – A vial of water that seeps from the oath‑breaker's boulder.
  ℎfill \textit{Drink it to remember a promise you broke. The memory is painful. The pain is the price – and the thread.} \texttt{[DEATH]}
ℹtem[\card{♦}{K}] \textbf{First Ancestor's Blessing} – A whispered promise from the Spirit‑Anchor Stone.
  ℎfill \textit{Once, you may reroll any failed roll involving a witnessed promise. The ancestor watches – and the spider weaves.} \texttt{[RITUAL]}
ℹtem[\card{♦}{A}] \textbf{The Ninth Warden's Coil} – A coil of silk said to belong to the legendary warden.
  ℎfill \textit{Use it to cut a bridge without angering the web‑law. The coil then unravels – and the spider's shadow falls on you.} \texttt{[LEGEND]}
\end{enumerate}

\section*{Quick use notes}
łabel{sec:lethaiar−quick−use}
\begin{itemize}
ℹtem Draw until you have all four suits: Spade = place, Heart = actor, Club = pressure, Diamond = leverage. Highest rank sets the main Timer (2–5 \(→\) 4, 6–10 \(→\) 6, J/Q/K \(→\) 8, A \(→\) 10).
ℹtem Diamonds are codified outcomes (cords, masks, shards) that change position rather than call for a roll.
ℹtem If any A appears, echo \textbf{valley‑omens} – silk cords humming in still air, the crimson river running clear, a bridge that sways without wind. The spider is near.
\end{itemize}

\subsection*{Additional Features (Cultural Mechanics)}
łabel{sec:lethaiar−mechanics}

\begin{itemize}
ℹtem \textbf{The Web‑Law:} The unwritten legal code of the Crimson Valley. Each warden carries a different clause in memory; together, they form a complete web. To cross a bridge, a traveler must answer a question about the web‑law. The answer changes with each warden. On a successful roll (Wits + Lore, DV 3), you may bypass the challenge; on a failure, you must offer a promise instead – and the spider will witness it.
ℹtem \textbf{Marks (Skin‑Tattoos):} Each Lethai−ar's Marks are a public record of promises witnessed. The Marks are cut from the skin after death and buried at the Ukwe's roots. To read a living Lethai−ar's Marks, one must know the phase of the crimson moon and the position of the valley's spirits. A character who studies a Lethai−ar's Marks for a scene may ask the GM one question about a promise that Lethai−ar has witnessed; the GM must answer truthfully – but the spider may also listen.
ℹtem \textbf{Spirit Ancestry:} Every Lethai−ar can name the specific spirit (mountain, river, grove) where their ancestor left the spirit world to become mortal. That spirit is their witness; they owe it a promise at birth, and that promise is never fully paid. Once per journey, a Lethai−ar may invoke their spirit ancestor to gain +1 die on a roll related to that spirit's domain (e.g., a river ancestor aids a crossing). The spider notes the favour.
ℹtem \textbf{Witness Economy:} Instead of tracking debt, the Crimson Valley uses \textbf{Witness} as currency. A character who witnesses a promise gains 1 Witness token. Tokens can be spent to: gain +1 die on a roll involving that promise, force an oath‑breaker to reroll a successful deception, or demand a favour from someone who witnessed the same promise. Witness tokens are lost when the promise is fulfilled or broken – and the spider watches the exchange.
ℹtem \textbf{Local Patrons (Syncretic Names):}
  \begin{itemize}
    ℹtem \textbf{The First Weaver (Inaea):} The spider who taught the first knot; mercy is a strand, not a cut. Her web is the valley's memory.
    ℹtem \textbf{The Serpent Judge (Isoka):} The fang that decides; a promise kept is a skin shed. She watches from the river stones.
    ℹtem \textbf{The Crimson Witness (The Witness):} The flame that sees through silk; no oath hidden from its light. The bridge glows when it passes.
    ℹtem \textbf{The Bridge Father (The Traveler):} The road that crosses the gorge; every crossing is a promise. He walks the strands at dusk.
  \end{itemize}
\end{itemize}

\subsection*{Lore Echoes (Cross‑Expansion Hooks)}
łabel{sec:lethaiar−lore}

\begin{itemize}
  ℹtem \textbf{Crimson Silk and the Way of Silk:} The silk of the Crimson Valley is prized across the continent for vow‑cords, well‑contracts, and bridge‑building. The Tkanki spider‑traders are the primary merchants; a single bolt can buy safe passage from Fhara to Dhahara. A Tkanki story‑seller may accept crimson silk as payment for a truth – but the spider will ask for a copy.
  ℹtem \textbf{The Ukwe's Roots and Sekogo:} The Ukwe's roots extend into the Crimson Valley. The Lethai−ar bury their dead Marks at the riverbank root, creating a unique relationship with the Speaking Tree. The Ukwe knows every promise ever made in the valley. A Sekogo druid may ask the Ukwe to witness a new oath, binding it with the weight of the valley – and the spider's thread.
  ℹtem \textbf{Web‑Law and the Covenant (Mykkiel):} The Covenant judges respect the web‑law as a valid legal system. A Lethai−ar truth‑cord is admissible in a Covenant court. Some judges have learned to read Marks. A Covenant justiciar may offer to translate a web‑law dispute into imperial law – but the translation always favours the spider.
  ℹtem \textbf{The Unmarked and the Hollow (Eastern Realms):} Lethai−ar who refuse all Marks become invisible to the web‑law – and to the Hollow. The Hollow cannot track them. Some Aelaerem seek out the Unmarked to learn the secret of blank skin. The Unmarked charge a promise: ``Do not seek me again.'' The spider does not visit them.
\end{itemize}

\subsection*{Lethai−ar Superstitions (burn for +1 die once)}
\begin{itemize}
  ℹtem \textbf{Bread to the Bridge:} Crumble a crust onto the first strand. Ignore the first \emph{Snapped Strand} penalty. \texttt{[SUPERSTITION]}
  ℹtem \textbf{Knot the Cord:} Tie a single knot in a truth‑cord. Cancel \emph{Web−Law Challenge} or take –1 die on your next negotiation (your choice). \texttt{[SUPERSTITION]}
  ℹtem \textbf{Name to the Root:} Whisper a dead warden's name into the Ukwe's root. Gain +1 Effect when crossing a bridge this scene, but mark Obligation to the First Ancestor – and the spider will know. \texttt{[SUPERSTITION]}
\end{itemize}

\subsection*{Thematic SB Spend Table}
łabel{sec:lethaiar−sb}

\paragraph{Minor Complications (1 SB)}
\begin{itemize}
ℹtem \textbf{Exposure:} Your actions draw unwanted attention from Silk‑Wardens or Mark‑Cutters.
ℹtem \textbf{Noise:} Sounds of your actions alert nearby bridge‑keepers or Tkanki traders.
ℹtem \textbf{Trace:} Evidence of your passage marks your route for Elders or the Unmarked.
ℹtem \textbf{Delay:} A brief but meaningful setback costs you time or safe crossing.
ℹtem \textbf{Supply Strain:} Mark +1 segment on a relevant resource timer.
\end{itemize}

\paragraph{Moderate Setbacks (2 SB)}
\begin{itemize}
ℹtem \textbf{Alarm Raised:} Local Silk‑Warden or Bridge‑Keeper becomes aware and begins responding.
ℹtem \textbf{Position Lost:} You lose advantageous ground/cover/stealth due to snapped strand or rising river.
ℹtem \textbf{Foe Appears:} An Oshiiran surveyor or Ashaani scout arrives on scene.
ℹtem \textbf{Gear Trouble:} A piece of equipment becomes Compromised/Neglected.
ℹtem \textbf{Lock/Barrier:} A simple obstacle now requires a test to overcome.
\end{itemize}

\paragraph{Serious Trouble (3 SB)}
\begin{itemize}
ℹtem \textbf{Reinforcements:} Additional Wardens, Mark‑Cutters, or bridge guards arrive.
ℹtem \textbf{Key Gear Breaks:} A crucial tool/weapon becomes temporarily unusable.
ℹtem \textbf{Major Twist:} The situation fundamentally changes – bridge cut / web‑law broken / First Bridge appears.
ℹtem \textbf{Rail Tick:} Advance a relevant campaign/front timer by 1 segment.
ℹtem \textbf{Condition Applied:} Mark \textbf{Fatigue 1/Harm 1/Condition} appropriate to fiction.
\end{itemize}

\paragraph{Major Turns (4+ SB)}
\begin{itemize}
ℹtem \textbf{Trap Springs:} A prepared danger activates with full effect.
ℹtem \textbf{Authority Arrival:} \textbf{Elder, Ninth Warden, or First Ancestor's Echo} intervenes.
ℹtem \textbf{Scene Shift:} The environment changes dramatically – \textbf{bridge collapses / river turns black / Ukwe root rises}.
ℹtem \textbf{Patron Omen:} Divine/arcane forces take notice – \textbf{omen appears / blessing lost / curse manifests}.
ℹtem \textbf{Narrative Pivot:} The story takes an unexpected turn that reframes objectives.
\end{itemize}

\paragraph{Region−Specific SB Options}
\begin{itemize}
ℹtem \textbf{Crimson Valley (Web‑Law):} Silk strands glow, truth‑cords tighten, the Ukwe's root hums a name.
ℹtem \textbf{Crimson Valley (Marks):} Tattoos itch, ochre stains spread, a dead Mark appears on living skin.
ℹtem \textbf{Crimson Valley (Bridges):} Ropes whisper, the gorge echoes with footsteps that are not there, a bridge sways toward the Hollow.
\end{itemize}

\subsection*{Pursuit Ladder (Near / Mid / Far)}
\begin{itemize}
ℹtem \textbf{Advance/Withdraw (Wits + Survival):} On a hit, shift one band. On a strong hit, also impose \emph{Swaying Bridge}.
ℹtem \textbf{Cut the Strand (Presence + Sway):} On a hit, freeze range for one exchange; if a truth‑cord is in your possession, +1 die.
ℹtem \textbf{Weather the Crossing (Resolve + Spirit):} On a hit in \emph{Fall from the Bridge}, you pick who catches the fallen; on a miss, the GM does.
\end{itemize}

\subsection*{Boss Seeds (Silk & Promise)}
\begin{itemize}
ℹtem \textbf{The Unweaver (Nine Broken Strands)} – \emph{Phases}: Snapped Cords \(→\) Orchestrated Falls \(→\) Bridge Coup. \emph{Cracks}: repair a broken strand, expose a false witness, survive a crossing without a promise.
ℹtem \textbf{The Silent Warden (Oath of the Unmarked)} – \emph{Phases}: Forgotten Promises \(→\) Procession of the Blank \(→\) Skin‑Baptism of an Elder. \emph{Cracks}: relic truth‑cord, survivor testimony, the Marks turning against them.
\end{itemize}

\begin{tcolorbox}[colback=black!3,colframe=black!40!white,title={Patronage & Power}]
Among the Lethai−ar, power flows through witnessed promises, silk strands, and the authority of the wardens. The web‑law is unwritten, but every warden carries a piece of it. Marks are public, stories are visible, and the Ukwe remembers everything. The spider – Inaea – does not rule, but she watches. Her threads are the bonds between families, the knots that hold the valley together, and the silent witness to every oath.

\textbf{For the GM:} Power in the Crimson Valley is measured in witnessed oaths and the trust of the bridge. Patronage often takes the form of truth‑cords, masks, or root shards – each tied to a specific promise. To emphasise this:
\begin{itemize}
ℹtem Tie rewards to objects that record or enforce promises (cords, masks, shards) that can be challenged, broken, or called in.
ℹtem Let rival wardens issue conflicting interpretations of the web‑law, forcing players to choose whose bridge to cross.
ℹtem Use the bridges, Mark‑cutter's studio, and the Ukwe's root as arenas for social contests, where a witnessed promise is a weapon and a broken strand is a death sentence.
ℹtem Remember that the spider is not a god – she is a presence. She does not speak, but the valley hums with her attention. When a promise is broken, the threads tremble.
\end{itemize}
In the Crimson Valley, the bridge remembers every step – and the spider forgets no promise.
\end{tcolorbox}

% Index entries
∈dex{Lethai−ar|see{Crimson Valley}}
∈dex{Theme generator}
∈dex{Crimson Valley generator}
∈dex{Places!Lethai−ar}
∈dex{People!Lethai−ar}
∈dex{Complications!Lethai−ar}
∈dex{Rewards!Lethai−ar}
∈dex{Silk−Warden}
∈dex{Mark−Cutter}
∈dex{Web−law}
∈dex{Truth−cord}
∈dex{Inaea (spider)}
∈dex{Ukwe (root)}
∈dex{Local Patrons!Lethai−ar}
